\section{Experimental Setup}
\label{sec:setup}

\paragraph{How large is the noise?}
We ran the baseline four times on ImageNet-R with one seed; the base task was
identical and only the SGD randomness of later tasks changed. The spread is
$\bar A=87.35\pm0.14$ and $A_B=82.59\pm0.23$; on ImageNet-A the spread of
$\bar A$ is $0.38$. A single run against a single run therefore cannot resolve
effects below roughly half a point. With the authors' code and partitions, our
runs reproduce all twelve numbers of the baseline's main table ($\bar A$ and
$A_B$ on six benchmarks) to within $0.01$.

\paragraph{Protocol.}
Four rules govern every comparison.
\textbf{(1) Pair everything:} effects are paired differences over seeds or
replicates, with 95\% $t$-intervals.
\textbf{(2) Same network for all heads:} a component that only changes
inference is evaluated, after each task, on the same test features as the
linear head. Each head draws randomness from a private generator and the
global random state is restored afterwards; we check that the linear head
reproduces the logged curve digit for digit. The standard deviation of a
paired head difference falls from $0.2$ to $0.03$--$0.05$ points.
Training-time components cannot use this rule and are compared with baseline
runs of the same seed and replicate.
\textbf{(3) Decide before looking:} each round had a written specification
fixing components, hyper-parameters, the decision rule and the minimum effect
before it ran; Table~\ref{tab:tests} (Appendix~\ref{app:tests}) lists every
test, including the failed ones.
\textbf{(4) A classical twin for every quantum component}
(Figure~\ref{fig:pipeline}): an RBF kernel on the same angles or the same
circuit without CNOTs for the PQK; MCR$^2$ for the Holevo quantity; a Gaussian
classifier for the POVM loss; ridge regression on the same statistics for the
PGM. Beating the baseline is not enough, since the baseline does not use the
information that the component and its twin share.

\paragraph{Read-only probe.}
To diagnose the separation regulariser we added a probe that measures,
without changing training, the cosine between the adapter gradients of the
classification, alignment and separation losses every 10 steps, and, after
each task on 20 training images per class, how well distances in each space
separate classes (AUC of different-class versus same-class pairs). With the
probe on, weights are bit-identical to a run without it.

\paragraph{Implementation.}
Training-time comparisons of QKSR use three replicates on ImageNet-R seed 1993
in one computing environment; the Holevo and residual-POVM objectives were
screened with one run each, against a pre-registered stopping rule. All runs
used an NVIDIA RTX PRO 6000 Blackwell GPU (up to three runs in parallel),
PyTorch 2.11 and timm 1.0.29.
Hyper-parameters are in Appendix~\ref{app:hyper}.
